%%% LaTeX-Vorlage Version 1.8 %%%

% Grundlegende Dokumenteneigenschaften gemäß DHBW-Vorgaben
\documentclass[a4paper,fontsize=11pt,oneside,parskip=half,headings=normal]{scrreprt} 
% \usepackage{showframe} % nur für Kontrolle der Ränder 

%%% Präambel einbinden (mit Festlegungen gemäß DHBW-Vorgaben) %%%
%%% Präambel %%%
% hier sollten keine Änderungen erforderlich sein
%
\usepackage[utf8]{inputenc}   % Zeichencodierung UTF-8 für Eingabe-Dateien
\usepackage[T1]{fontenc}      % Darstellung von Umlauten im PDF

\usepackage{listings}         % für Einbindung von Code-Listings
\lstset{numbers=left,numberstyle=\tiny,numbersep=5pt,texcl=true}
\lstset{literate=             % erlaubt Sonderzeichen in Code-Listings 
{Ö}{{\"O}}1
{Ä}{{\"A}}1
{Ü}{{\"U}}1
{ß}{{\ss}}2
{ü}{{\"u}}1
{ä}{{\"a}}1
{ö}{{\"o}}1
{€}{{\euro}}1
}

\usepackage[
  inner=35mm,outer=15mm,top=25mm,
  bottom=20mm,foot=12mm,includefoot
]{geometry}                 % Einstellungen für Ränder

\usepackage[ngerman]{babel} % Spracheinstellungen Deutsch
\usepackage[babel,german=quotes]{csquotes} % deutsche Anf.zeichen
\usepackage{enumerate}      % anpassbare Nummerier./Aufz.
\usepackage{graphicx}       % Einbinden von Grafiken
\usepackage[onehalfspacing]{setspace} % anderthalbzeilig

\usepackage{blindtext}      % Textgenerierung für Testzwecke
\usepackage{color}          % Verwendung von Farbe 

\usepackage{acronym}        % für ein Abkürzungsverzeichnis

\usepackage[                % Biblatex
  backend=biber,
  bibstyle=_dhbw_authoryear,maxbibnames=99,
  citestyle=authoryear,     
  uniquename=true, useprefix=true,
  bibencoding=utf8]{biblatex}
%kein Punkt am Ende bei \footcite
%http://www.golatex.de/footcite-ohne-punkt-am-schluss-t4865.html
\renewcommand{\bibfootnotewrapper}[1]{\bibsentence#1}


%Reihenfolge der Autorennamen
%   
% http://golatex.de/viewtopic,p,80448.html#80448
% Argumente: siehe http://texwelt.de/blog/modifizieren-eines-biblatex-stils/
\DeclareNameFormat{sortname}{% Bibliographie
  \ifnum\value{uniquename}=0 % Normalfall
    \ifuseprefix%
      {%
         \usebibmacro{name:family-given}
           {\namepartfamily}
           {\namepartgiveni}
           {\namepartprefix}
           {\namepartsuffixi}%
       }
      {%
         \usebibmacro{name:family-given}
           {\namepartfamily}
           {\namepartgiveni}
           {\namepartprefixi}
           {\namepartsuffixi}%
       }%
  \fi
  \ifnum\value{uniquename}=1% falls nicht eindeutig, abgek. Vorname 
      {%
         \usebibmacro{name:family-given}
           {\namepartfamily}
           {\namepartgiveni}
           {\namepartprefix}
           {\namepartsuffix}%
       }%
  \fi
  \ifnum\value{uniquename}=2% falls nicht eindeutig, ganzer Vorname 
      {%
         \usebibmacro{name:family-given}
           {\namepartfamily}
           {\namepartgiven}
           {\namepartprefix}
           {\namepartsuffix}%
       }%
  \fi   
  \usebibmacro{name:andothers}}

\DeclareNameFormat{labelname}{% für Zitate
  \ifnum\value{uniquename}=0 % Normalfall
    \ifuseprefix%
      {%
         \usebibmacro{name:family-given}
           {\namepartfamily}
           {\empty}
           {\namepartprefix}
           {\namepartsuffixi}%
       }
      {%
         \usebibmacro{name:family-given}
           {\namepartfamily}
           {\empty}
           {\namepartprefixi}
           {\namepartsuffixi}%
       }%
  \fi
  \ifnum\value{uniquename}=1% falls nicht eindeutig, abgek. Vorname 
      {%
         \usebibmacro{name:family-given}
           {\namepartfamily}
           {\namepartgiveni}
           {\namepartprefix}
           {\namepartsuffix}%
       }%
  \fi
  \ifnum\value{uniquename}=2% falls nicht eindeutig, ganzer Vorname 
      {%
         \usebibmacro{name:family-given}
           {\namepartfamily}
           {\namepartgiven}
           {\namepartprefix}
           {\namepartsuffix}%
       }%
  \fi   
  \usebibmacro{name:andothers}}
      
  
\DeclareFieldFormat{extrayear}{% = the 'a' in 'Jones 1995a'
  \iffieldnums{labelyear}
    {\mknumalph{#1}}
    {\mknumalph{#1}}}        

\renewcommand*{\multinamedelim}{\addslash}
\renewcommand*{\finalnamedelim}{\addslash}
\renewcommand*{\multilistdelim}{\addslash}
\renewcommand*{\finallistdelim}{\addslash}

\renewcommand{\nameyeardelim}{~}

% Literaturverzeichnis: Doppelpunkt zwischen Name (Jahr): Rest 
% http://de.comp.text.tex.narkive.com/Tn1HUIXB/biblatex-authoryear-und-doppelpunkt
\renewcommand{\labelnamepunct}{\addcolon\addspace}

% damit die Darstellung für Vollzitate von Primärquellen in 
% Fußnoten später auf "nicht fett" geändert werden kann 
% (nur für Zitate von Sekundärliteratur relevant)
\newcommand{\textfett}[1]{\textbf{#1}}

% für Zitate von Sekundärliteratur:
\newcommand{\footcitePrimaerSekundaer}[4]{%
  \renewcommand{\textfett}[1]{##1}%
  \footnote{\fullcite[#2]{#1}, zitiert nach \cite[#4]{#3}}%  
  \renewcommand{\textfett}[1]{\textbf{##1}}%
}

% Im Literaturverzeichnis: Autor (Jahr) fett
\renewbibmacro*{author}{%
  \ifboolexpr{%
    test \ifuseauthor%
    and
    not test {\ifnameundef{author}}
  }
    {\usebibmacro{bbx:dashcheck}
       {\bibnamedash}
       {\usebibmacro{bbx:savehash}%
        \textfett{\printnames{author}}%
        \iffieldundef{authortype}
          {\setunit{\addspace}}
          {\setunit{\addcomma\space}}}%
     \iffieldundef{authortype}
       {}
       {\usebibmacro{authorstrg}%
        \setunit{\addspace}}}%
    {\global\undef\bbx@lasthash
     \usebibmacro{labeltitle}%
     \setunit*{\addspace}}%
  \textfett{\usebibmacro{date+extrayear}}}

% Sonderfall: Quelle ohne Autor, aber mit Herausgeber
% Name des Herausgebers wird fett gedruckt
\renewbibmacro*{bbx:editor}[1]{%
  \ifboolexpr{%
    test \ifuseeditor%
    and
    not test {\ifnameundef{editor}}
  }
    {\usebibmacro{bbx:dashcheck}
       {\bibnamedash}
       {\textfett{\printnames{editor}}%
        \setunit{\addcomma\space}%
        \usebibmacro{bbx:savehash}}%
     \usebibmacro{#1}%
     \clearname{editor}%
     \setunit{\addspace}}%
    {\global\undef\bbx@lasthash
     \usebibmacro{labeltitle}%
     \setunit*{\addspace}}%
  \textfett{\usebibmacro{date+extrayear}}}

% Anpassungen für deutsche Sprache
\DefineBibliographyStrings{ngerman}{%
	nodate = {{o.J.}},
	urlseen = {{Abruf:}},
	ibidem = {{ebenda}}
}

% keine Anführungszeichen beim Titel im Literaturverzeichnis
\DeclareFieldFormat[article,book,inbook,inproceedings,manual,misc,phdthesis,thesis,online,report]{title}{#1\isdot}

\newcommand{\literaturverzeichnis}{%
% nur Literaturverzeichnis
% (als eigenes Kapitel)
\phantomsection
\addcontentsline{toc}{chapter}{Literaturverzeichnis}
\spezialkopfzeile{Literaturverzeichnis}
\defbibheading{lit}{\chapter*{Literaturverzeichnis}}
\label{chapter:quellen}
\printbibliography[heading=lit,notkeyword=ausblenden]
} % mit DHBW-spezifischen Einstellungen

\usepackage{hyperref}       % URL-Formatierung, klickbare Verweise

\usepackage{tocloft}        % für Verzeichnis der Anhänge

\usepackage{multirow}

\newcounter{anhcnt}
\setcounter{anhcnt}{0}
\newlistof{anhang}{app}{}

\newcommand{\anhang}[1]{%
  \refstepcounter{anhcnt}
  \setcounter{anhteilcnt}{0}
  \section*{Anhang \theanhcnt: #1}
  \addcontentsline{app}{section}{\protect\numberline{Anhang \theanhcnt}#1}\par
}

\newcounter{anhteilcnt}
\setcounter{anhteilcnt}{0}

\newcommand{\anhangteil}[1]{%
	\refstepcounter{anhteilcnt}
	\subsection*{Anhang~\arabic{anhcnt}/\arabic{anhteilcnt}: #1}
	\addcontentsline{app}{subsection}{\protect\numberline{Anhang \theanhcnt/\arabic{anhteilcnt}}#1}\par
}

\renewcommand{\theanhteilcnt}{Anhang \theanhcnt/\arabic{anhteilcnt}}

% vgl. S. 4 Paket-Beschreibung tocloft 	
% Einrückungen für Anhangverzeichnis
\makeatletter
\newcommand{\abstaendeanhangverzeichnis}{
\renewcommand*{\l@section}{\@dottedtocline{1}{0em}{5.5em}}
\renewcommand*{\l@subsection}{\@dottedtocline{2}{2.3em}{6.5em}}
}
\makeatother

% Abbildungs- und Tabellenverzeichnis
% Bezeichnungen
\renewcaptionname{ngerman}{\figurename}{Abb.}
\renewcaptionname{ngerman}{\tablename}{Tab.}
% Einrückungen
\makeatletter
\renewcommand*{\l@figure}{\@dottedtocline{1}{0em}{2.3em}}
\renewcommand*{\l@table}{\@dottedtocline{1}{0em}{2.3em}}
\makeatother


\usepackage{chngcntr}                % fortlaufende Zähler für Fußnoten, Abbildungen und Tabellen
\counterwithout{figure}{chapter}
\counterwithout{table}{chapter}
\counterwithout{footnote}{chapter}

\usepackage[automark]{scrlayer-scrpage} 
%% Definitionen für Kopf- und Fußzeile auf normalen Seiten
\defpagestyle{kopfzeile}
{% Kopfdefinition
  (\textwidth,0pt)    % Länge der oberen Linie,Dicke der oberen Linie       
  {} % Definition für linke Seiten im doppelseitigen Layout
  {} % Definition für rechte Seiten im doppelseitigen Layout      
  {  % Definition für Seiten im einseitigen Layout
	\makebox[0pt][l]{\rightmark}% 
	\makebox[\linewidth]{}% 
  }        
  (\textwidth, 0.4pt) % Untere Linienlänge, Untere Liniendicke
}
{% Fußdefinition
  (\textwidth,0pt)    % Obere Linienlänge, Obere Liniendicke
  {} % Definition für linke Seiten im doppelseitigen Layout
  {} % Definition für rechte Seiten im doppelseitigen Layout
  {  % Definition für Seiten im einseitigen Layout
    \makebox[\linewidth]{}%
    \makebox[0pt][r]{\pagemark}%
  }
  (\textwidth, 0pt)   % Länge der unteren Linie,Dicke der unteren Linie
}

%% Definitionen für Kopf- und Fußzeile auf ersten Seiten eines Kapitels
\defpagestyle{kapitelkopfzeile}
{% Kopfdefinition
  (\textwidth,0pt)    % Länge der oberen Linie,Dicke der oberen Linie       
  {} % Definition für linke Seiten im doppelseitigen Layout
  {} % Definition für rechte Seiten im doppelseitigen Layout      
  {}  % Definition für Seiten im einseitigen Layout
  (\textwidth, 0pt) % Untere Linienlänge, Untere Liniendicke
}
{% Fußdefinition
  (\textwidth,0pt)    % Obere Linienlänge, Obere Liniendicke
  {} % Definition für linke Seiten im doppelseitigen Layout
  {} % Definition für rechte Seiten im doppelseitigen Layout
  {  % Definition für Seiten im einseitigen Layout
    \makebox[\linewidth]{}%
    \makebox[0pt][r]{\pagemark}%
  }
  (\textwidth, 0pt)   % Länge der unteren Linie,Dicke der unteren Linie
}

%% Definitionen für Kopf- und Fußzeile im Anhang und bei Quellenverzeichnisse
\newcommand{\spezialkopfzeileBezeichnung}{}
\defpagestyle{spezialkopfzeile}
{% Kopfdefinition
  (\textwidth,0pt)    % Länge der oberen Linie,Dicke der oberen Linie       
  {} % Definition für linke Seiten im doppelseitigen Layout
  {} % Definition für rechte Seiten im doppelseitigen Layout      
  {  % Definition für Seiten im einseitigen Layout
	\makebox[0pt][l]{\spezialkopfzeileBezeichnung}% 
	\makebox[\linewidth]{}% 
  }        
  (\textwidth, 0.4pt) % Untere Linienlänge, Untere Liniendicke
}
{% Fußdefinition
  (\textwidth,0pt)    % Obere Linienlänge, Obere Liniendicke
  {} % Definition für linke Seiten im doppelseitigen Layout
  {} % Definition für rechte Seiten im doppelseitigen Layout
  {  % Definition für Seiten im einseitigen Layout
    \makebox[\linewidth]{}%
    \makebox[0pt][r]{\pagemark}%
  }
  (\textwidth, 0pt)   % Länge der unteren Linie,Dicke der unteren Linie
}
            
\newcommand\spezialkopfzeile[1]{%
  \renewcommand\spezialkopfzeileBezeichnung{#1}
  \pagestyle{spezialkopfzeile}
}
                
% Standard-Pagestyle auswählen
\pagestyle{kopfzeile}

% keine Kopfzeile anzeigen auf Seiten, auf denen ein 
% Kapitel beginnt oder das Inhalts-/Abbildungs-/Tabellenverzeichnis steht 
\renewcommand{\chapterpagestyle}{kapitelkopfzeile}
\tocloftpagestyle{kapitelkopfzeile}

		 % für schöne Kopfzeilen 

\usepackage{textcomp}            % erlaubt EUR-Zeichen in Eingabedatei
\usepackage{eurosym}             % offizielles EUR-Symbol in Ausgabe
\renewcommand{\texteuro}{\euro}  % ACHTUNG: nach hyperref aufrufen!

\usepackage{scrhack}             % stellt Kompatibilität zw. KOMA-Script
                                 % (scrreprt) und anderen Paketen her
                                 
% Anpassung der Abstände bei Kapitelüberschriften
% (betrifft auch Inhalts-, Abbildungs- und Tabellenverzeichnis)
\renewcommand*\chapterheadstartvskip{\vspace*{-\topskip}}
\newcommand{\myBeforeTitleSkip}{1mm}
\newcommand{\myAfterTitleSkip}{10mm}
\setlength\cftbeforetoctitleskip{\myBeforeTitleSkip}
\setlength\cftbeforeloftitleskip{\myBeforeTitleSkip}
\setlength\cftbeforelottitleskip{\myBeforeTitleSkip}

\setlength\cftaftertoctitleskip{\myAfterTitleSkip}
\setlength\cftafterloftitleskip{\myAfterTitleSkip}
\setlength\cftafterlottitleskip{\myAfterTitleSkip}                                                            
%%% Ende der Präambel %%%


%%% Name der eigenen Literatur-Datenbank (ggf. anpassen) %%%
\bibliography{includes/chatbot.bib}

\begin{document}
%%% Deckblatt einbinden %%% 
% Anpassungen nötig (Name, Titel etc.)
% HIER EDITIEREN: 
% Typ der Arbeit (für Deckblatt und ehrenwörtliche Erklärung)
% - bitte Zutreffendes auswählen
\newcommand{\typMeinerArbeit}{1. Projektarbeit} 
%\newcommand{\typMeinerArbeit}{2. Projektarbeit} 
%\newcommand{\typMeinerArbeit}{Seminararbeit} 
% \newcommand{\typMeinerArbeit}{Bachelorarbeit} 

% Thema der Arbeit (für ehrenwörtliche Erklärung, ohne Umbrüche)
% HIER EDITIEREN: 
\newcommand{\themaMeinerArbeit}{Was ist ein Chatbot?}

% Vorname, Name der Autorin/des Autors (für Titelseite und Metadaten)
% HIER EDITIEREN:
\newcommand{\meinName}{Max Weiberle}

\thispagestyle{empty}

\begin{spacing}{1}
\begin{center}	
~\vspace{0mm}

% HIER EDITIEREN: Titel der Arbeit
{\sffamily
\LARGE  
% \Large  % bei sehr langen Titeln ggf. etwas kleinere Schriftart wählen
\textbf{Was ist ein Chatbot?}

\bigskip
\textbf{}
}


\vspace{15mm}

% Typ wird automatisch eingefügt (oben festlegen)
{\Large \typMeinerArbeit}

\vspace{1cm}

% HIER ggf. EDITIEREN
vorgelegt am 18. Mai 2022

\vspace{15mm}

Fakultät Wirtschaft
\medskip

Studiengang Wirtschaftsinformatik
\medskip

% HIER EDITIEREN: Kurs eintragen
Kurs WWI2021G

\vspace{10mm}

von

\vspace{10mm}

% Vorname und Name wird automatisch eingefügt (oben festlegen) 
{\large\textsc{\meinName}}

\vspace{10mm}
\end{center}

\vfill

% HIER EDITIEREN: Name des Unternehmens, Name der Betreuerin/des Betreuers
\begin{tabular}{ll}
Betreuer in der Ausbildungsstätte: & DHBW Stuttgart: \\
\hspace{0.4\linewidth} & \\
Komm.ONE & Prof., Dr., Manfred Sander \\
Stefan Fuchs \\
Studierendenbetreuer \\
\\
Unterschrift der Betreuerin/des Betreuers \\
\end{tabular}


\vspace{1cm}
%(etwas Platz für die Unterschrift der Betreuerin/des Betreuers aus der Ausbildungsstätte)
\end{spacing}

% falls ein Vertraulichkeitsvermerk erforderlich ist,
% die Kommentarzeichen in den nachfolgenden Zeilen entfernen:
 
%\begin{center}
%\small
%\textbf{Vertraulichkeitsvermerk}:
%Der Inhalt dieser Arbeit darf weder als Ganzes noch in Auszügen \\
%Personen außerhalb des Prüfungs- und Evaluationsverfahrens zugänglich gemacht werden, sofern keine anders lautende Genehmigung des Dualen Partners vorliegt. 
%\end{center}

% Meta-Daten für PDF-Datei basierend auf obigen Angaben
\hypersetup{pdftitle={\themaMeinerArbeit}}
\hypersetup{pdfauthor={\meinName}}
\hypersetup{pdfsubject={\typMeinerArbeit\ DHBW Stuttgart \the\year}}


%%% Umstellung der Seiten-Nummerierung auf i, ii, iii ... %%%
\pagenumbering{Roman} 

\cleardoublepage

%%% Inhalts-, Abbildungs-, Tabellenverzeichnisse %%%
% sollen einzeilig gesetzt werden, um Platz zu sparen 
\begin{spacing}{1}
\tableofcontents

\clearpage
\thispagestyle{kapitelkopfzeile}
\listoffigures
\phantomsection
\addcontentsline{toc}{chapter}{Abbildungsverzeichnis} % Abb.verz. ins Inh.verz. aufnehmen

\end{spacing}

%%% Umstellung der Seiten-Nummerierung auf 1, 2, 3 ... %%%
\cleardoublepage
\pagenumbering{arabic}

%%% Ihr eigentlicher Inhalt %%%
% Empfehlung: strukturieren Sie Ihren Text in einzelnen Dateien 
% und binden Sie diese hier mit \input{includes/dateiname.tex} ein

\chapter{Einleitung} % (fold)
\label{cha:Einleitung}

Durch die Corona-Pandemie und die damit verbundene Entwicklung möglichst viel im Homeoffice zu arbeiten wurde die Verwendung von VPN-Verbindungen für uns bereits zum Alltag. Zudem kommt die sich ankündigende Energiekrise. Da ein Anstieg der Energiekosten um 50 Prozent aktuell als realistisch betrachtet wird und von 30.000 Menschen gesprochen wird, die noch dieses Jahr arbeitslos werden könnten, scheint es wichtiger denn je zu sein Energiekosten zu sparen. \footcite[Vgl.][]{Gries.2022} Energieexperten berufen sich auf Studien, die eine Einsparung von bis zu fünf Prozent durch Homeoffice vorhersagen. \footcite[Vgl.][S. 1 f.]{Thaler.2022} Des Weiteren haben sich Homeoffice und flexible Arbeitszeiten zu starken Kriterien für die Berufswahl entwickelt. Besonders junge Menschen legen großen Wert darauf. Millenials würden Umfragen zu Folge doppelt so häufig den Arbeitsplatz wechseln als Babyboomer, wenn sie mit der zeitlichen oder örtlichen Flexibilität nicht zufrieden sind. \footcite[Vgl.][S. 1 f.]{Hahn.2021}

Für die Firma Komm.One, als regionales Dienstleistungsunternehmen, die das Rechenzentrum von Baden-Württemberg betreut und für einen möglichst reibungslosen Ablauf in den Rathäusern und Kommunen des Landes sorgt, ist es essentiell eine zuverlässige und beständige Arbeit zu leisten. Das ist nur möglich, wenn die Belegschaft durchgängig arbeiten kann. In den letzten zwei Jahren und vor allem zu den Corona-Hochzeiten konnte dies an den Standorten nicht gewährleistet werden, da bei der Infektion einer Person alle Kontaktperson in Quarantäne mussten. Als IT Firma blieb also die Möglichkeit des Homeoffice als sinnvolle Alternative. In dieser Zeit hat sich gezeigt, dass Homeoffice sowohl Vor- als auch Nachteile gegenüber dem zuvor standardmäßigen Arbeiten vor Ort hat. Je nach Aufgabe und Verantwortung ist dies natürlich unterschiedlich. Für die Komm.One hat sich das Homeoffice als wertvolle und auch in Zukunft interessante Alternative bestätigt. Aktuell haben allein am Standort Reutlingen etwa 180 von 250 Mitarbeitern einen Homeoffice-Vertrag, der bis zu fünf Homeoffice-Tage ermöglicht.

Um diese Menge an Homeoffice möglich zu machen, ist bei der Komm.One eine VPN-Infrastruktur gewachsen, die aus historischen Gründen, auf die in der Arbeit genauer eingegangen wird, nicht optimal ist und überarbeitet werden muss. Es werden nur für interne Mitarbeiter aktuell drei verschiedene VPN-Lösungen verwendet. Deshalb verfolgt die Komm.One das Ziel einer Harmonisierung der internen VPN-Struktur durchzuführen. In dieser Arbeit sollen die möglichen VPN-Lösungen evaluiert und die für die Komm.One beste Lösung gefunden werden.

% chapter (end)

\chapter{Firewalls und VPN} % (fold)
\label{cha:Firewalls und VPN}

Um die Rolle, die VPN in nahezu allen Unternehmen spielt verstehen zu können, müssen erst ein paar grundlegende Informationen zu Netzwerken erläutert und verstanden werden.

\section{Firewalls} % (fold)
\label{sec:Firewalls}

Grundsätzlich ist eine Firewall alles, Hardware, Software oder eine Kombination aus beidem, was die Übertragung von digitalen Paketen filtern kann. Dem liegen zwei Sicherheitsfunktionen zu Grunde:

\begin{description}
	\item Filtern von Paketen: Digitale Pakete, die durch die Firewall in das zu schützende Netzwerk hinein, aber auch aus dem zu schützenden Netzwerk hinaus wollen werden auf Basis der erstellten Sicherheitsregeln gefiltert.
	\item Anwendungsproxy: Schützt individuelle Computer und ermöglicht dennoch das Nutzen von Netzwerkdiensten. Hierfür wird der IP-Informationsfluss unterbrochen. \footcite[Vgl.][S. 99]{WHITMAN.2018}
\end{description}

\begin{figure}[htb]
	\centering
	\includegraphics[width=9cm]{graphics/dhbw.png}
	\caption[Idee eines Firewall-Systems]{Idee eines Firewall-Systems. \footnotemark}
	\label{abb:Firewall}
\end{figure}
\footnotetext{Entnommen aus: \cite[S. 357]{Pohlmann.2022}}

In Abbildung \ref{abb:Firewall} ist die erste Funktion gut zu erkennen. Die Firewall lässt die grüne Verbindung durch, weil sie nach den Sicherheitsrichtlinien in Ordnung ist, während die orangene Linie abgeblockt wird. Ein schöner Vergleich ist hier der Pförtner oder Sicherheitspersonal am Eingang. Sie kümmern sich darum, dass die bestehenden Regeln eingehalten werden. Alle Personen müssen durch dafür vorgesehene Bereiche das Gebäude betreten oder verlassen, werden auf ihre Zugangsberechtigung geprüft und müssen durch einen Metalldetektor gehen und potenziell gefährliche Dinge abgeben. Falls es zu einem Sicherheitsfehler kommt und eine Person die Kontrolle umgeht muss eine Warnung ausgegeben werden und ein anderes Team sich um das Problem kümmern. Analog dazu funktioniert auch die Firewall, welche kontrolliert über welche Protokolle und Dienste zugegriffen wird und wer mit wem kommunizieren darf. All diese Daten werden protokolliert. \footnotemark
\footnotetext{Vgl. \cite[S. 100]{WHITMAN.2018} und \cite[S. 358]{Pohlmann.2022}}

Folgende Aufgaben haben Firewall-Systeme im Allgemeinen:

\begin{itemize}
	\item Einschränken von Zugriffen von außerhalb des Netzwerks: Die offensichtlichste Gefahr ist der Angriff von außerhalb auf das eigene Netzwerk. Aus diesem Grund muss eine Firewall jedes einzelne Paket auf die individuell festgelegten Regeln überprüfen.
	\item Einschränken von unautorisierten Zugriffen von innerhalb des Netzwerks: In manchen Fällen ist es einfacher ein Netzwerk vor äußeren Einflüssen zu schützen als vor den Gefahren, die von unaufmerksamen Arbeitnehmern ausgehen können. Besonders zu beachten sind hier beispielsweise mobile Datenträger, die infizierte Dateien in das System einschleusen können. Besonders häufig kommen auch E-Mail Angriffe mit infizierten, oder ausführbaren Dateien, die beim Herunterladen oder Ausführen das gesamte Netzwerk betreffen können.
	\item Zugangskontrolle auf Nutzerebene: Es wird überprüft, welche Nutzer über das Firewall-System mit welchen Systemen kommunizieren dürfen.
	\item Zugangskontrolle auf Datenebene: Es wird überprüft, welche Daten eines definierten Nutzers und Systems über das Firewall-System übertragen werden dürfen.
	\item Rechteverwaltung: Es wird festgelegt, mit welchen Protokollen und Diensten und zu welchen Zeiten über das Firewall-System kommuniziert werden darf.
	\item Kontrolle auf Anwendungsebene: Es wird überprüft, ob Aktivitäten stattfinden, die nicht zu Anwendungen oder definierten Aufgabenstellungen gehören, generell unerwünscht sind oder Schaden am System verursachen können. Beispiele hierfür wären Spam oder Malware.
	\item Entkopplung von Diensten: Dienste werden vom Netzwerk entkoppelt, um zu verhindern, dass Fehler oder Sicherheitslücken der einzelnen Dienste zu erfolgreichen Angriffen führen können.
	\item Beweissicherung und Protokollauswertung: Jegliche Verbindungen und sicherheitsrelevante Ereignisse werden protokolliert. Diese können für die Erkennung von Sicherheitsverletzungen und zur Beweissicherung von Handlungen der Nutzer ausgewertet werden.
	\item Alarmierung: Besondere sicherheitsrelevante Ereignisse werden dem Security-Management gemeldet, so dass im Ernstfall schnell gehandelt werden kann.
	\item Verbergen der internen Netzstruktur: Im Idealfall ist aus dem unsicheren Netz heraus nicht zu erkennen, wie das zu schützende Netz strukturiert ist oder wie viele Systeme sich in diesem Netz befinden. \footnotemark
\end{itemize}
\footnotetext{Vgl. \cite{WHITMAN.2018} S. 103 ff. und \cite{Pohlmann.2022} S. 359 f.}

Allerdings ist es bei Weitem nicht mit einer Firewall getan. Um die Sicherheit eines Netzwerkes so hoch wie möglich zu halten, müssen mehrere Sicherheitssysteme zusammenarbeiten.

% section Firewalls (end)

\section{VPN} % (fold)
\label{sec:VPN}

Ein weiterer Schritt um das zu schützende Netz sicherer zu machen, aber vor allem auch um ganz neue Möglichkeit zu eröffnen ist das Einrichten von VPN-Verbindungen. Diese Technologie ermöglicht nicht nur einen Zusammenschluss von zwei Netzwerken, die an verschiedenen geografischen Standorten existieren. Zunehmend gewinnt auch die Verwendung im Bereich von Homeoffice immer mehr an Relevanz. VPN ermöglicht es mehrere Netzwerke, wie in Abbildung \ref{abb:VPN-Modell} dargestellt, oder Geräte zu einem Netzwerk zusammenzufassen, ohne dass sie nebeneinander sein müssen oder durch ein Kabel verbunden werden. In den allermeisten Fällen wird das Internet als Verbindung zwischen den Netzwerken oder Geräten verwendet. Hier wird die Verbindung über ein unsicheres Netz mit einem Tunnel dargestellt. Um diese Verbindungen brauchbar und sicher zu machen, ist eine Verschlüsselung essentiell. \footcite[Vgl.][S. 261 ff.]{WHITMAN.2018}

\begin{figure}[htb]
	\centering
	\includegraphics[width=12cm]{graphics/dhbw.png}
	\caption[Vereinfachtes Modell eines VPN-Aufbaus]{Vereinfachtes Modell eines VPN-Aufbaus \footnotemark}
	\label{abb:VPN-Modell}
\end{figure}
\footnotetext{Entnommen aus: \cite{WHITMAN.2018} S. 296}

Die Verschlüsselung kann bildlich wie ein Tunnel verstanden werden und eine VPN-Verbindung wird häufig auch als VPN-Tunnel bezeichnet. Die beiden Router mit IPSec zeigen genau diese Verschlüsselung. IPSec ist eine in der IT weit verbreitete Verschlüsselungsmethode und der aktuelle Standard. \footnotemark
\footnotetext{Vgl. \cite{o.V..2020} und \cite{Pohlmann.2022} S. 403 ff. auch im Folgenden}

\begin{figure}[htb]
	\centering
	\includegraphics[width=10cm]{graphics/dhbw.png}
	\caption[Realisierungsformen von IPSec]{Realisierungsformen von IPSec \footnotemark}
	\label{abb:IPSec}
\end{figure}
\footnotetext{Entnommen aus: \cite{Pohlmann.2022} S. 411}

Für einen Verbindungsaufbau werden, wie in Abbildung \ref{abb:IPSec} dargestellt, entweder zwei IPSec-Clients, zwei IPSec-Gateways oder jeweils ein Gateway und ein Client benötigt.

% section VPN (end)

% chapter (end)

\chapter{Relevanz} % (fold)
\label{cha:Relevanz}

\section{Vorteile gegenüber herkömmlichen Systemen} % (fold)
\label{sec:Vorteile gegenüber herkömmlichen Systemen}

In der nachfolgenden Abbildung (Abb \ref{abb:VorteileChatbotsPolen})

\begin{figure}[htb]
  \centering
  \includegraphics[width=12cm]{graphics/dhbw.png}
  \caption[Vorteile für Firmen bei Verwendung von Chatbots im Kundenservice für Unternehmen in Polen]{Vorteile für Firmen bei Verwendung von Chatbots im Kundenservice für Unternehmen in Polen. \footnotemark}
  \label{abb:VorteileChatbotsPolen}
\end{figure}
\footnotetext{Entnommen aus: \cite{Bisser.2021}}

Werden die Vorteile der Verwendung von Chatbots aufgezeigt und deren Wirksamkeit. \\
Als die zwei größten Vorteile zeigt die Abbildung die eingesparte Zeit beim Kontaktieren von Kunden und die Möglichkeit Kunden rund um die Uhr betreuen zu können.

% section Vorteile gegenüber herkömmlichen Systemen (end)

\section{Nachteile gegenüber herkömmlichen Systemen} % (fold)
\label{sec:Nachteile gegenüber herkömmlichen Systemen}

In der nachfolgenden Abbildung (Abb \ref{abb:NachteileChatbots})

\begin{figure}[htb]
  \centering
  \includegraphics[width=12cm]{graphics/dhbw.png}
  \caption[Nachteile bei der Verwendung eines Chatbots aus Sicht der Kunden]{Nachteile bei der Verwendung eines Chatbots aus Sicht der Kunden. \footnotemark}
  \label{abb:NachteileChatbots}
\end{figure}
\footnotetext{Entnommen aus: \cite{.2021}}

Werden, die laut Kunden, größten Nachteile von Chatbots aufgelistet und nach Häufigkeit sortiert.

% section Nachteile gegenüber herkömmlichen Systemen (end)

% chapter (end)


%%% Ende des eigentlichen Inhalts %%%


%%% Quellenverzeichnisse (keine Anpassung nötig) %%%
\clearpage
\literaturverzeichnis
%%% Ende Quellenverzeichnisse %%%


%%% Erklärung (keine Anpassungen nötig) %%%
% steht ganz am Ende des Dokuments
\cleardoublepage
\clearpage

\thispagestyle{empty}

{\LARGE\textsf{\textbf{Erklärung}}\bigskip}

% \typMeinerArbeit und \themaMeinerArbeit werden in deckblatt.tex definiert
Ich versichere hiermit, dass ich meine \typMeinerArbeit\ mit dem Thema: \emph{\themaMeinerArbeit} selbstständig verfasst und keine anderen als die angegebenen Quellen und Hilfsmittel benutzt habe.
Ich versichere zudem, dass die eingereichte elektronische Fassung mit der gedruckten Fassung übereinstimmt.

\vspace{3cm}

\begin{center}
\begin{tabular}{ccc}
(Ort, Datum) & \hspace{0.3\linewidth} & (Unterschrift)
\end{tabular}
\end{center}
\end{document}
